\TieudeT{PHONG TỎA VẬT LÍ}
\subsubsection*{PHẦN I. Thí sinh trả lời câu hỏi từ câu 1 đến câu 18. Mỗi câu hỏi thí sinh chỉ chọn một phương án}
\setcounter{ex}{0}
\begin{ex}%Tailieuchuan
Khi một vật đứng yên và \textbf{không} có lực tác dụng, vật sẽ
\choice
{\True tiếp tục đứng yên}
{bắt đầu chuyển động}
{chuyển động nhanh dần}
{chuyển động chậm dần}
\loigiai{}
\end{ex}
\begin{ex}%book
Khi nói về tác dụng của lực lên vật, phát biểu nào dưới đây đúng?
\choice
{\True Khi vật thay đổi vận tốc thì bắt buộc phải có lực tác dụng vào vật}
{Vật bắt buộc phải chuyển động theo hướng của lực tác dụng vào nó}
{Nếu không còn lực nào tác dụng vào vật đang chuyển động thì vật phải lập tức dừng lại}
{Một vật không thể liên tục chuyển động mãi mãi nếu không có lực nào tác dụng vào nó}
\loigiai{
Theo định luật 1 Newton A đúng, B,C,D sai.
} \haidong{8} \end{ex}
\begin{ex}%book %[1P2Y2-1][Loai1]
Phát biểu nào sau đây là đúng?
\choice
{Một vật sẽ đứng yên nếu không chịu tác dụng của lực nào}
{Một vật chịu tác dụng của hai lực cân bằng thì đứng yên}
{Một vật chỉ ở trạng thái cân bằng khi vật đứng yên}
{\True Hai lực cân bằng tác dụng vào một vật đang đứng yên thì vật đứng yên}
\loigiai{
Phát biểu hai lực cân bằng tác dụng vào một vật đang đứng yên thì vật đứng yên là đúng.
}
\haidong{8} \end{ex}
\begin{ex}%book%[1P2B3-1][Loai1]
Khi nói về tác dụng của lực lên vật, phát biểu nào sau đây là đúng?
\choice
{Nếu không còn lực nào tác dụng vào vật đang chuyển động thì vật phải lập tức dừng lại}
{Một vật không thể liên tục chuyển động mãi mãi nếu không có lực nào tác dụng vào nó}
{Vật bắt buộc phải chuyển động theo hướng của lực tác dụng vào nó}
{\True Khi vật thay đổi vận tốc thì bắt buộc phải có lực tác dụng vào vật}
\loigiai{
\begin{itemize}
\item Lực có hai tác dụng là gây biến dạng vật và gây biến đổi chuyển động của vật. 
\item Khi vật thay đổi vận tốc, tức là có sự biến đổi chuyển động thì bắt buộc phải có lực tác dụng vào vật.
\end{itemize}
}
\haidong{8} \end{ex}
\begin{ex}%book
Theo định luật 1 Newton thì
\choice
{lực là nguyên nhân duy trì chuyển động}
{\True một vật sẽ giữ nguyên trạng thái đứng yên hoặc chuyển động thẳng đều nếu nó không chịu tác dụng của lực nào}
{một vật không thể chuyển động được nếu hợp lực tác dụng lên nó bằng 0}
{mọi vật đang chuyển động đều có xu hướng dừng lại do quán tính}
\loigiai{
Áp dụng định luật I Newton suy ra B đúng.
} \haidong{8} \end{ex}
\begin{ex}%book
Hợp lực tác dụng vào một vật đang chuyển động thẳng đều bằng hợp lực tác dụng vào vật 
\choice
{chuyển động tròn đều}
{tự do}
{chuyển động nhanh dần đều}
{\True đứng yên}
\loigiai{
Hợp lực tác dụng vào một vật đang chuyển động thẳng đều bằng hợp lực tác dụng vào vật đứng yên.
} \haidong{8} \end{ex}
\begin{ex}%book%[1P2Y2-1][Loai1] 
Nôi dung nào dưới đây là của định luật 1 Newton?
\choice
{\True Nếu một vật không chịu tác dụng của lực nào hoặc chịu tác dụng của hợp lực có độ lớn bằng không, thì vật đang đứng yên sẽ tiếp tục đứng yên, đang chuyển động sẽ tiếp tục chuyển động thẳng đều}
{Trong mọi trường hợp, khi vật A tác dụng lên vật B một lực, thì vật B cũng tác dụng lại vật A một lực, hai lực này có độ lớn bằng nhau}
{Gia tốc của một vật cùng hướng với lực tác dụng lên vật, có độ lớn tỉ lệ thuận với độ lớn của lực và tỉ lệ nghịch với khối lượng của vật}
{Khối lượng là đại lượng đặc trưng cho mức quán tính của vật}
\loigiai{
Nội dung định luật I: Nếu một vật không chịu tác dụng của lực nào hoặc chịu tác dụng của hợp lực có độ lớn bằng không, thì vật đang đứng yên sẽ tiếp tục đứng yên, đang chuyển động sẽ tiếp tục chuyển động thẳng đều.
}
\haidong{8} \end{ex}
\begin{ex}%book
Nếu một vật đang chuyển động mà tất cả các lực tác dụng vào nó bỗng nhiên ngừng tác dụng thì
\choice
{vật lập tức dừng lại}
{vật chuyển động chậm dần rồi dừng lại}
{vật chuyển động chậm dần trong một thời gian, sau đó sẽ chuyển động thẳng đều}
{\True vật chuyển động thẳng đều}
\loigiai{
Nếu một vật đang chuyển động mà tất cả các lực tác dụng vào nó bỗng nhiên ngừng tác dụng thì vật chuyển động thẳng đều.
} \haidong{8} \end{ex}
\begin{ex}%book
Một vật nằm yên trên mặt bàn là do
\choice
{vật chỉ chịu tác dụng của lực hút Trái Đất}
{không có lực tác dụng lên vật}
{các lực tác dụng lên vật có cường độ quá nhỏ}
{\True lực hút của Trái Đất lên vật cân bằng với phản lực của bàn}
\loigiai{
Vật nằm yên trên mặt bàn là do lực hút của Trái đất lên vật cân bằng với phản lực của bàn.
} \haidong{8} \end{ex}
\begin{ex}%book %[1P2B2-1][Loai1]
Người ta kéo một vật chuyển động trên mặt phẳng nằm ngang không ma sát. Nếu thôi tác dụng lực lên vật thì vật sẽ
\choice
{dừng lại ngay}
{tiếp tục chuyển động chậm dần rồi dừng lại}
{\True chuyển ngay sang trạng thái chuyển động thẳng đều}
{thay đổi vận tốc}
\loigiai{ 
Do vật chuyển động trên mặt phẳng nằm ngang không ma sát nên khi thôi tác dụng lực thì vật sẽ chuyển động thẳng đều.
}
\haidong{8} \end{ex}
\begin{ex}%book
Khi đang đi xe đạp trên đường nằm ngang, nếu ta ngừng đạp, xe vẫn đi tiếp chứ chưa dừng ngay. Đó là nhờ
\choice
{trọng lượng của xe}
{lực ma sát}
{\True quán tính của xe}
{phản lực của mặt đường}
\loigiai{
Theo quán tính mọi vật bảo toàn trạng thái đứng yên, chuyển động của vật.
Nên xe đạp nếu đang đạp mà dừng thì theo quán tính sẽ đi tiếp một đoạn rồi mới dừng lại.
} \haidong{8} \end{ex}
\begin{ex}%book
Vật nào sau đây chuyển động theo quán tính?
\choice
{Vật chuyển động tròn đều}
{Vật chuyển động trên quỹ đạo thẳng}
{\True Vật chuyển động thẳng đều}
{Vật chuyển động rơi tự do}
\loigiai{
Theo định luật I Newton khi vật không chịu tác dụng của lực nào hoặc của các lực có hợp lực bằng không, thì vật đang đứng yên sẽ tiếp tục đứng yên, đang chuyển động sẽ chuyển động thẳng đều do vậy C đúng.

} \haidong{8} \end{ex}
\begin{ex}%book
Trường hợp nào sau đây nhất định có liên quan tới quán tính
\choice
{Chiếc bè trôi trên sông}
{Vật rơi trong không khí}
{\True Dũ quần áo cho sạch bụi}
{Vật rơi tự do}
\loigiai{
} \haidong{8} \end{ex}
\begin{ex}%book
Hiện tượng nào sau đây \textbf{không} thể hiện tính quán tính?
\choice
{Khi bút máy bị tắc mực người ta vẫy mực để mực văng ra}
{\True Viên bi có khối lượng lớn lăn xuống máng nghiêng nhanh hơn viên bi có khối lượng nhỏ}
{Ôtô đang chuyển động thì tắt máy nó vẫn chạy thêm 1 đoạn nữa rồi mới dừng lại}
{Một người đứng trên xe buýt, xe hãm phanh đột ngột, người có xu hướng ngã về phía trước}
\loigiai{
Ta thấy viên bi có khối lượng lớn lăn xuống máng nghiêng nhanh hơn viên bi có khối lượng nhỏ không thể hiện tính quán tính.
} \haidong{8} \end{ex}
\begin{ex}%book%[1P2Y2-1][Loai1] 
Khi nói về quán tính của một vật, phát biểu nào sau đây \textbf{không} đúng?
\choice
{Nguyên nhân duy trì chuyển động khi các lực tác dụng lên vật mất đi chính là tính quán tính của vật}
{Khi hợp lực tác dụng lên vật bằng không, vật sẽ đứng yên hoặc chuyển động thẳng đều theo quán tính}
{\True Những vật có khối lượng rất nhỏ thì không có quán tính}
{Quán tính là tính chất của mọi vật bảo toàn vận tốc của mình khi không chịu tác dụng của lực}
\loigiai{
Mọi vật đều có quán tính.
}
\haidong{8} \end{ex}
\begin{ex}%book%[1P2B2-1][Loai1]
Hành khách ngồi trên xe ôtô đang chuyển động, xe bất ngờ rẽ sang phải. Theo quán tính hành khách sẽ
\choice
{nghiêng sang phải}
{\True nghiêng sang trái}
{ngả người về phía sau}
{chúi người về phía trước}
\haidong{8} \end{ex}
\begin{ex}%book
Điều gì sẽ có thể xảy ra với người lái xe máy chạy ngay sau một xe tải nếu xe tải dừng đột ngột?
\choice
{Không có vấn đề gì}
{Xe máy sẽ bị văng về phía sau}
{Xe máy sẽ bị đổ sang phải}
{\True Xe máy sẽ bị đâm vào xe tải}
\loigiai{
Do phản xạ của người lái xe máy là không tức thời mà cần có một khoảng thời gian dù rất ngắn để nhận ra xe tải đã dừng và ấn chân vào phanh.

Do xe có quán tính nên dù đã hãm cũng không thể dừng lại ngay mà cần có thời gian để dừng hẳn. Trong hai khoảng thời gian trên nếu xe máy đi hết khoảng cách giữa hai xe và có thể đâm vào xe tải.
} \haidong{8} %<MyLT1>
\end{ex}
\begin{ex}%book
Cho các phát biểu sau:
\begin{itemize}
\item[(1)] Định luật 1 Newton còn được gọi là định luật quán tính.
\item[(2)] Mọi vật đều có xu hướng bảo toàn vận tốc của mình.
\item[(3)] Chuyển động thẳng đều được gọi là chuyển động theo quán tính.
\item[(4)] Quán tính là tính chất của mọi vật có xu hướng bảo toàn vận tốc cả về hướng và độ lớn.
\end{itemize}
Số phát biểu đúng là	
\choice
{1}
{\True 4}
{3}
{2}
\loigiai{
Số phát biểu đúng là 4 .
} \haidong{8} \end{ex}















\subsubsection*{PHẦN 2. Thí sinh trả lời từ câu 1 đến câu 4. Trong mỗi ý a), b), c), d) ở mỗi câu, thí sinh chọn đúng hoặc sai.}
\setcounter{ex}{0}
\begin{ex}%Tailieuchuan
Một quả bóng được đặt trên mặt đất và sau đó bị đẩy bằng tay để nó lăn trên mặt phẳng nằm ngang.
\choiceTF[t]
{\True Khi quả bóng lăn, quán tính của nó làm cho quả bóng tiếp tục di chuyển theo phương ngang dù tay không còn đẩy nó} 
{Khi quả bóng lăn trên mặt đất, các lực cản không ảnh hưởng đến chuyển động của nó} 
{\True Nếu không có lực tác dụng nào nữa, quả bóng sẽ tiếp tục lăn với vận tốc không đổi mãi mãi} 
{\True Khi quả bóng lăn và lực đẩy từ tay không còn, quả bóng sẽ dần dừng lại do các lực cản tác dụng lên nó} 
\loigiai{}
\end{ex}

\begin{ex}%Tailieuchuan
Một người đang đứng trên một chiếc xe buýt đang chạy thẳng với tốc độ không đổi. Đột nhiên, xe buýt phanh gấp.
\choiceTF[t]
{\True Người đứng trên xe buýt sẽ bị ngã về phía trước khi xe buýt phanh gấp vì cơ thể họ có quán tính muốn giữ nguyên trạng thái chuyển động ban đầu} 
{Nếu người đó đang cầm một chiếc cốc, chiếc cốc sẽ bị rơi về phía sau khi xe buýt phanh gấp vì lực quán tính tác dụng lên cốc theo hướng ngược lại} 
{Người đứng trên xe buýt không bị ngã về phía trước khi xe phanh gấp vì chân của họ đã bám chặt vào sàn xe} 
{\True Khi xe buýt phanh gấp, tất cả các vật dụng trong xe buýt sẽ bị đẩy về phía trước do quán tính của chúng} 
\loigiai{}
\end{ex}
\begin{ex}%Tailieuchuan
Khi bạn phơi quần áo vừa giặt xong, bạn thường rũ mạnh quần áo trước khi treo lên dây phơi.
\choiceTF[t]
{\True Khi rũ áo quần mạnh, nước bên trong quần áo bị văng ra ngoài vì quán tính làm cho nước tiếp tục chuyển động khi áo quần đột ngột thay đổi hướng} 
{\True Khi rũ áo quần, các sợi vải cũng được làm thẳng ra do lực tác động của tay, giúp quần áo khô nhanh hơn và ít nhăn hơn} 
{Khi rũ áo quần mạnh, nước trong quần áo sẽ tự động bay hơi nhanh hơn vì tác động của không khí xung quanh} 
{Khi rũ áo quần mạnh, nước sẽ không bị văng ra ngoài nếu quần áo vẫn còn quá ướt, vì nước bám chặt vào vải hơn} 
\loigiai{}
\end{ex}
\begin{ex}%Tailieuchuan
Một người đang di chuyển bằng xe máy với vận tốc không đổi trên đường thẳng.
\choiceTF[t]
{Nếu người lái xe máy ngừng tăng ga, xe máy sẽ ngay lập tức dừng lại do không còn lực tác dụng lên nó} 
{\True Nếu người lái xe máy ngừng tăng ga, xe máy sẽ tiếp tục di chuyển một khoảng trước khi dừng lại do quán tính} 
{\True Nếu xe máy di chuyển với vận tốc không đổi, lực kéo của động cơ và lực cản tác dụng lên xe là cân bằng nhau} 
{Nếu xe máy đi lên dốc, nó sẽ tiếp tục di chuyển với vận tốc không đổi mà không cần thêm lực kéo từ động cơ} 
\loigiai{}
\end{ex}


\newpage 

\TieudeT{PHONG TỎA VẬT LÍ}
\subsubsection*{PHẦN I. Thí sinh trả lời câu hỏi từ câu 1 đến câu 18. Mỗi câu hỏi thí sinh chỉ chọn một phương án}
\setcounter{ex}{0}
\begin{ex}%book%[1P2Y2-1][Loai1]
Khi nói về tác dụng của lực đối với chuyển động, phát biểu nào sau đây là đúng?
\choice
{Một vật bất kì chịu tác dụng của một lực có độ lớn tăng dần thì chuyển động nhanh dần}
{Nếu không có lực tác dụng vào vật thì vật không thể chuyển động}
{\True Một vật có thể chịu tác dụng đồng thời của nhiều lực mà vẫn chuyển động thẳng đều}
{Không vật nào có thể chuyển động ngược chiều với lực tác dụng lên nó}
\loigiai{
Theo định luật I Niu - tơn ta có khi một vật chịu tác dụng của nhiều lực mà tổng hợp của các lực đó bằng 0 thì vật đang chuyển động sẽ chuyển động thẳng đều.
}
\haidong{8} \end{ex}
\begin{ex}%Tailieuchuan
Khi một quả bóng đang lăn trên mặt đất và dừng lại, điều này chứng tỏ
\choice
{không có lực nào tác dụng lên quả bóng}
{\True có lực cản tác dụng lên quả bóng}
{quả bóng mất quán tính}
{quả bóng thay đổi hướng chuyển động}
\loigiai{

}
\end{ex}
\begin{ex}%book%[1P2B2-1][Loai1]
Phát biểu nào sau đây chưa chính xác?
\choice
{Nếu chỉ có một lực duy nhất tác dụng lên vật thì vận tốc của vật thay đổi}
{\True Nếu có lực tác dụng lên vật thì độ lớn vận tốc của vật bị thay đổi}
{Nếu có nhiều lực tác dụng lên vật mà các lực này cân bằng nhau thì vận tốc của vật không thay đổi}
{Nếu vận tốc của vật không đổi thì không có lực nào tác dụng lên vật hoặc các lực tác dụng lên vật cân bằng nhau}
\loigiai{
Lực còn có thể làm vật bị biến dạng.
}
\haidong{8} \end{ex}
\begin{ex}%book%[1P2B3-1][Loai1]
Phát biểu nào sau đây là \textbf{sai}? Một vật chịu tác dụng của một lực khi
\choice
{\True vật đó đứng yên}
{vật đó thay đổi hình dạng}
{vật đó thay đổi hướng chuyển động}
{vật đó chuyển động nhanh lên hay chậm đi}
\haidong{8} \end{ex}
\begin{ex}%book
Một vật sẽ đứng yên hoặc chuyển động thẳng đều khi
\choice
{\True các lực tác dụng vào vật cân bằng nhau}
{các lực tác dụng vào vật cùng chiều với nhau}
{các lực tác dụng vào vật ngược chiều nhau, độ lớn khác nhau}
{các lực tác dụng vào vật vuông góc với nhau}
\loigiai{
Theo định luật I Newton thì một vật sẽ đứng yên hoặc chuyển động thẳng đều khi các lực tác dụng vào vật cân bằng nhau
} \haidong{8} %<MyLT1>
\end{ex}
\begin{ex}%book
Định luật 1 Newton còn được gọi là
\choice
{\True Định luật quán tính}
{Định luật li tâm}
{Định luật phi quán tính}
{Định luật hướng tâm}
\loigiai{
} \haidong{8} \end{ex}
\begin{ex}%book
Quán tính là tính chất mọi vật có xu hướng bảo toàn
\choice
{tốc độ của mình}
{\True vận tốc của mình}
{khối lượng của mình}
{gia tốc của mình}
\loigiai{
Do có quán tính mà mọi vật có xu hướng bảo toàn vận tốc cả về hướng và độ lớn.
} \haidong{8} \end{ex}
\begin{ex}%book
Một xe ô tô đang chuyển động thẳng với vận tốc không đổi là $20\ m/s$. Hợp lực tác dụng lên ô tô có độ lớn bằng
\choice
{$20\ N$}
{\True $0\ N$}
{$10\ N$}
{$-20\ N$}
\loigiai{
Vì ô tô chuyển động thẳng đều nên theo định luật I Newton ô tô chịu tác dụng của hợp lực bằng 0 .
} \haidong{8} \end{ex}
\begin{ex}%book %[1P2B2-1][Loai1]
Một vật chuyển động thẳng với vận tốc $\vec v$ không đổi, tác dụng vào vật nhiều lực có phương khác nhau sao cho hợp lực của chúng bằng 0, thì chuyển động của vật sẽ như thế nào?
\choice
{Vật dừng lại ngay}
{Vật chuyển động chậm dần đều với vận tốc ban đầu là v}
{Vật chuyển động nhanh dần đều với vận tốc ban đầu là v}
{\True Vật chuyển động thẳng đều với vận tốc là v}
\loigiai{
Vật chuyển động thẳng đều với vận tốc là v.
}
\haidong{8} \end{ex}
\begin{ex}%book
Khi đang đi xe máy trên đường nằm ngang, nếu ta phanh xe, xe vẫn tự di chuyển. Đó là nhờ
\choice
{trọng lượng của xe}
{lực ma sát nhỏ}
{\True quán tính của xe}
{phản lực của mặt đường}
\loigiai{
Do quán tính xe vẫn đi tiếp một đoạn trước khi dừng.
} \haidong{8} %<MyLT1>
\end{ex}
\begin{ex}%book
Khi đang đi ô tô trên đường nằm ngang, nếu ta phanh gấp, xe vẫn tự đi chuyển. Đó là nhờ
\choice
{trọng lượng của xe}
{lực ma sát nhỏ}
{\True quán tính của xe}
{phản lực của mặt đường}
\loigiai{
Quán tính là tính chất bảo toàn trạng thái đứng yên hoặc chuyển động của vật. do vậy khi xe đang đi mà phanh gấp thì xe vẫn dịch lên một đoạn trước khi dừng.
} \haidong{8} \end{ex}
\begin{ex}%book
Vật nào sau đây chuyển động theo quán tính?
\choice
{Vật chuyển động tròn đều}
{Vật chuyển động trên một đường thẳng}
{Vật rơi tự do từ trên cao xuống không ma sát}
{\True Vật chuyển động khi tất cả các lực tác dụng lên vật mất đi}
\loigiai{
Theo định luật I Newton vật chuyển động theo quán tính khi tất cả các lực tác dụng lên vật bị mất đi.
} \haidong{8} \end{ex}
\begin{ex}%book
Lúc chạy để tránh con chó sói đuổi bắt, con cáo thường thoát thân bằng cách bất thình lình rẽ ngoặt sang hướng khác, đúng vào lúc con chó sói định ngoạm cắn nó. Cáo làm vậy là vì
\choice
{\True theo quán tính, chó sói sẽ chạy theo hướng cũ một đoạn nữa nên cáo thoát được}
{theo quán tính, cả sói và cáo đều chạy theo hướng cũ một đoạn nữa}
{cáo phản xạ theo tự nhiên để đánh lạc hướng sói}
{cáo theo quán tính làm lạc hướng sói}
\loigiai{
Theo quán tính, chó sói sẽ chạy theo hướng cũ một đoạn nữa nên cáo thoát được.
} \haidong{8} \end{ex}
\begin{ex}%book %[1P2B2-1][Loai1]
Đặt một cốc đầy nước lên trên tờ giấy học trò. Tác dụng rất nhanh một lực F theo phương nằm ngang của tờ giấy thì hiện tượng gì sẽ xảy ra với tờ giấy và cốc nước?
\choice
{\True Tờ giấy rời khỏi cốc nước mà nước vẫn không đổ}
{Tờ giấy chuyển động và cốc nước chuyển động theo}
{Tờ giấy chuyển động về một hướng, cốc nước chuyển động theo hướng ngược lại}
{Tờ giấy bị đứt ở chỗ đặt cốc nước}
\loigiai{ 
Hiện tượng xảy ra là tờ giấy rời khỏi cốc nước mà nước vẫn không đổ. Do khi tác dụng lực trong thời gian ngắn do quán tính chiếc cốc không kịp thay đổi vận tốc tức là vận tốc vẫn giữ nguyên (bằng 0).
}
\haidong{8} \end{ex}
\begin{ex}%book %[1P2Y2-1][Loai1]
Khi một xe buýt tăng tốc đột ngột thì các hành khách thì 
\choice
{dừng lại ngay}
{\True ngả người về phía sau}
{chúi người về phía trước}
{ngả người sang bên cạnh}
\loigiai{ 
Khi một xe buýt tăng tốc đột ngột thì các hành khách ngả người về phía sau do quán tính.
}
\haidong{8} \end{ex}
\begin{ex}%book
Một hành khách ngồi ở cuối xe phàn nàn rằng, do lái xe phanh gấp mà một túi sách ở phía trước bay về phía anh ta làm anh ta bị đau. Người này nói đúng hay sai? Vì sao?
\choice
{Nói đúng vì theo quán tính túi bị bay về phía sau}
{Nói đúng vì túi bị bay về phía trước}
{\True Nói sai vì theo quán tính túi bị bay về phía trước}
{Nói sai vì các lực lúc này triệt tiêt túi không di chuyển}
\loigiai{
Anh ta nói sai vì theo quán tính vật bảo toàn trạng thái đứng yên hay chuyển động của nó do vậy khi xe phanh gấp chiếc túi phải bay về phía trước.
} \haidong{8} \end{ex}
\begin{ex}%book
Có hai phát biểu sau đây:
\begin{itemize}
\item[(1)] Do có quán tính, máy bay không thể tức thời đạt tới tốc độ đủ lớn để cất cánh. Nó phải tăng tốc dần trên đường băng mới cất cánh được. Khi hạ cánh, nó đang có tốc độ lớn nên phải hãm dần trên đường băng mới dừng lại được.
\item[(2)] Khi xe đang chạy nhanh mà dừng đột ngột, người ngồi trên xe sẽ bị xô về phía trước, có thể bị lao khỏi ghế hoặc bị chấn thương do va chạm mạnh vào bộ phận của xe phía trước chỗ ngồi của mình. Dây an toàn có tác dụng giữ cho người khỏi lao về phía trước khi xe dừng đột ngột.
\end{itemize}
Phát biểu nào là đúng?
\choice
{(1) đúng, (2) sai}
{\True (1) đúng, (2) đúng}
{(1) sai, (2) sai}
{(1) sai, (2) đúng}
\loigiai{
Theo quán tính cả (1) và (2) đều đúng.
} \haidong{8} \end{ex}
\begin{ex}%book
Có hai phát biểu sau đây:
\begin{itemize}
\item[(1)] Một vật đang đứng yên. Ta có thể kết luận, vật không chịu tác dụng của lực nào.
\item[(2)] Một hành khác ngồi ở cuối xa. Nếu lái xe phanh gấp thì một túi sách ở phía trước bay về phía anh ta. 
\end{itemize}
Phát biểu đúng là
\choice
{(1) đúng, (2) sai}
{(1) đúng, (2) đúng}
{\True (1) sai, (2) sai}
{(1) sai, (2) đúng}
\loigiai{
(1) sai vì theo định luật ${I}$ Newton vật đang đứng yên có thể do các lực tác dụng có hợp lực bằng không.
(2) sai vì theo quán tính túi sách phải bay vể phía trước.
} \haidong{8} \end{ex}















\subsubsection*{PHẦN 2. Thí sinh trả lời từ câu 1 đến câu 4. Trong mỗi ý a), b), c), d) ở mỗi câu, thí sinh chọn đúng hoặc sai.}
\setcounter{ex}{0}
\begin{ex}%Tailieuchuan
Một chiếc xe đạp đang di chuyển trên đường thẳng và bằng phẳng với tốc độ không đổi.
\choiceTF[t]
{Khi người đi xe đạp ngừng đạp, xe sẽ ngay lập tức dừng lại do không còn lực tác dụng lên nó} 
{\True Khi người đi xe đạp ngừng đạp, xe sẽ tiếp tục di chuyển thêm một đoạn trước khi dừng lại do quán tính} 
{\True Khi xe đạp di chuyển với vận tốc không đổi, lực đẩy của người đi xe và lực cản tác dụng lên xe là cân bằng nhau} 
{Nếu đường đi lên dốc, xe đạp sẽ tiếp tục di chuyển với vận tốc không đổi mà không cần thêm lực đẩy từ người đi xe} 
\loigiai{}
\end{ex}
\begin{ex}%Tailieuchuan
Một chiếc xe ô tô đang di chuyển trên một con đường thẳng và bất ngờ phanh gấp.
\choiceTF[t]
{\True Khi xe phanh gấp, hành khách ngồi trong xe bị đổ về phía trước do quán tính của cơ thể muốn giữ nguyên trạng thái chuyển động} 
{Khi xe phanh gấp, hành khách ngồi trong xe không bị ảnh hưởng vì lực phanh chỉ tác dụng lên bánh xe, không tác dụng lên người trong xe} 
{\True Nếu hành khách ngồi trong xe không thắt dây an toàn, họ có thể bị đẩy về phía trước mạnh hơn do không có lực nào giữ họ lại khi xe dừng đột ngột} 
{Khi xe phanh gấp, tất cả các vật trong xe đều giữ nguyên vị trí vì xe và các vật đã dừng lại cùng một lúc} 
\loigiai{}
\end{ex}
\begin{ex}%Tailieuchuan
Một chiếc xe buýt đang di chuyển với tốc độ không đổi trên đường thẳng. Một hành khách đứng trong xe buýt không nắm vào bất cứ thứ gì.
\choiceTF[t]
{\True Khi xe buýt di chuyển với tốc độ không đổi, hành khách trong xe sẽ cảm thấy ổn định và không bị đẩy về phía trước hoặc phía sau} 
{\True Nếu xe buýt đột ngột giảm tốc, hành khách sẽ bị đẩy về phía trước} 
{Nếu xe buýt đột ngột tăng tốc, hành khách sẽ không cảm thấy sự thay đổi và vẫn đứng yên tại chỗ} 
{Nếu xe buýt di chuyến với tốc độ không đổi và rẽ trái, hành khách sẽ cảm thấy bị đẩy về phía trước} 
\loigiai{}
\end{ex}
\begin{ex}%DoanhNamDinh
\immini[thm]{Một giáo viên yêu cầu học sinh tiến hành thí nghiệm gồm một hòn bi coi là chất điểm và hai máng nghiêng $(1)$ và $(2)$: Thả hòn bi ở độ cao $h$ so với mặt phẳng ngang từ máng nghiêng $(1)$ và sau đó lăn lên máng nghiêng $(2)$.}{
\includegraphics[width=5cm]{images/2025_06_04_af03f1d88c157fe9bf98g-126(1)}}
\choiceTF[t]
{\True Nếu bỏ qua mọi ma sát và lực cản thì hòn bi lăn lên máng nghiêng $(2)$ tới độ cao bằng độ cao $h$}
{\True Hòn bi lăn không tới độ cao bằng độ cao $h$ là vì có ma sát giữa hòn bi và máng nghiêng}
{\True Khi hạ thấp dần độ cao của máng nghiêng $(2)$, hòn bi lăn trên máng $(2)$ được một đoạn dài hơn so với với lần trước đó}
{\True Nếu không có ma sát giữa hòn bi và máng nghiêng và nếu máng nghiêng $(2)$ nằm ngang và dài vô hạn thì hòn bi sẽ lăn mãi mãi với vận tốc không đổi}
\loigiai{
\begin{itemchoice}
\itemch Trong điều kiện lí tưởng, hòn bi lăn lên máng nghiêng (2) tới độ cao bằng độ cao h do không có ma sát giữa hòn bi và máng nghiêng.
$\Rightarrow$ Đúng
\itemch Hòn bi lăn không tới độ cao bằng độ cao h là vì có ma sát giữa hòn bi và máng nghiêng.
$\Rightarrow$ Đúng
\itemch Khi hạ thấp dần độ cao của máng nghiêng (2), hòn bi lăn trên máng (2) được một đoạn dài hơn so với với lần trước đó.
$\Rightarrow$ Đúng
\itemch Nếu không có ma sát giữa hòn bi và máng nghiêng và nếu máng nghiêng (2) nằm ngang rất dài thì hòn bi sẽ lăn mãi mãi với vận tốc không đổi theo định luật I Newton.
$\Rightarrow$ Đúng
\end{itemchoice}
}
\end{ex}














